\documentclass[colorback,accentcolor=tud1c,linedtoc,bibtotoc]{tudreport}
\usepackage{fontspec}	%use LuaLaTeX or XeLaLaTeX as LaTeX engine
\usepackage{ngerman}
\usepackage[ngerman]{hyperref}
\usepackage{caption}
\usepackage{subcaption}
\usepackage{graphicx}
\usepackage{placeins}
\usepackage{amsmath}
\usepackage[stable]{footmisc}
\usepackage{ziffer}
\usepackage{upgreek}
\usepackage{parskip}

\hypersetup{
  pdftitle={TU Darmstadt - Physikalisches Praktikum für Fortgeschrittene},
  pdfauthor={Author1, Author2},
  pdfsubject={Versuch Nummer},
  pdfview=FitH,
}

%To get the fonts working, just download the .ttf files from the TU Corporate Design webpage
%https://www.tu-darmstadt.de/kommunikation_und_medien/corporate_design_1/index.de.jsp
%and upload them to the project, then uncomment the lines below

% \setmainfont{Charter}[ 
% Extension = .ttf,
% UprightFont = *-Regular,
% BoldFont = *-Bold,
% ItalicFont = *-Italic]

% \setmonofont{Stafford}[ 
% Extension = .ttf,
% UprightFont = *-Regular]

% \setsansfont{FrontPage}[ 
% Extension = .ttf,
% UprightFont = *-Medium,
% SmallCapsFont = *-Regular]

%BEGINN TITELSEITE

\title{Versuchsname}

\subtitle{Author1  \\ Author2 
}

\subsubtitle{\textsc{Betreuer:  \hfill Versuchsdatum: 
}}

\author{Author1, Author2}

\institution{Physikalisches Praktikum \\für Fortgeschrittene \\ Abteilung X \\ Institut für XXXphysik}

%ENDE TITELSEITE

%ANFANG DOKUMENT
\begin{document}




%Titelseite einfügen
\maketitle

\paragraph{Erklärung zur Selbständigkeit:} \ \\
Hiermit versichern wir die vorliegende Ausarbeitung des Fortgeschrittenen-Praktikums ohne Hilfe Dritter nur mit den angegebenen Quellen und Hilfsmitteln angefertigt zu haben. Alle Stellen, die aus Quellen entnommen wurden, sind als solche kenntlich gemacht. Diese Arbeit hat in gleicher oder ähnlicher Form noch keiner Prüfungsbehörde vorgelegen.

\begin{center} \
	\linebreak[1]\linebreak[1]\linebreak[1]\linebreak[1]
	\rule{5cm}{0.1mm}\\
	Author1
	\linebreak[1]\linebreak[1]\linebreak[1]\linebreak[1]\linebreak[1]
	\rule{5cm}{0.1mm}\\
	Author2\\
\end{center}

%Inhaltsverzeichnis einfügen
\tableofcontents

%ANFANG INHALT
\chapter*{Ziel des Versuchs}

\chapter{Grundlagen}
\section{1. Thema}
\subsection{1. Unterthema}

\chapter{Fazit}
\newpage
\begin{thebibliography}{99}

\end{thebibliography}
  
\newpage
  
\chapter*{Anhang}
\addcontentsline{toc}{chapter}{Anhang} 
\end{document}
